\fichatitulo{21 · RECUPERAÇÃO}{Preprint v1 · 2021}{BEIR}
{\small Thakur et al. \textit{{BEIR}: A Heterogeneous Benchmark for Zero-shot Evaluation of Information Retrieval Models}. Preprint v1 · 2021. \href{https://arxiv.org/pdf/2104.08663v1}{Fonte consultada}.\par}

\campo{Problema e objetivo}
Os recuperadores generalizam para domínios sem ajuste específico?

\campo{Principal contribuição · síntese interpretativa}
Propõe um benchmark heterogêneo para avaliar recuperação fora da distribuição de treinamento.

\campo{Método / natureza da evidência}
Compara abordagens lexicais, densas e de reordenação em tarefas variadas, sob transferência sem treinamento no domínio de destino.

\campo{Resultado ou argumento principal}
Na versão consultada, BM25 constitui uma referência robusta; reordenadores obtêm resultados fortes com maior custo computacional.

\campo{Excertos-chave · seleção literal}
\begin{itemize}
\excerto{1}{BM25 is a robust baseline}{Resumo.}
\excerto{2}{high computational costs}{Resumo.}
\end{itemize}

\campo{Limitações e análise crítica}
Análise da ficha: médias entre coleções ocultam diferenças por domínio. A evidência histórica não determina qual recuperador será melhor em discursos parlamentares em português.

\campo{Aproveitamento na dissertação · proposta}
Metodologia: manter baseline lexical e comparar qualidade, custo e desempenho por tipo de consulta.

\registro{Fonte consultada nesta preparação: Resumo e apresentação do benchmark. Versão arXiv v1: 17 conjuntos e nove modelos; não confundir com versões posteriores. Chave: \texttt{thakurEtAl2021beir}.}
